\chapter{Opérations sur les limites et suites géométriques}
\section*{Exercice 10.1 -- Terme dominant}
En divisant numérateur et dénominateur par $n^2$, on obtient $(5+2/n)/(1-3/n^2)\to5$. Le dénominateur est non nul pour les grands indices et sa limite vaut $1$.

\section*{Exercice 10.2 -- Conjugué}
Multiplier par le conjugué donne $4n/(\sqrt{n^2+4n}+n)=4/(\sqrt{1+4/n}+1)\to2$. L'expression initiale était une forme indéterminée.

\section*{Exercice 10.3 -- Géométriques}
Les deux premières ont module $0{,}7^n\to0$, donc convergent vers zéro. Pour la troisième, les termes pairs tendent vers $+\infty$ et les impairs vers $-\infty$ ; elle n'a ni limite réelle ni limite infinie unique.

\section*{Exercice 10.4 -- Suite affine}
Le point fixe est $6$. Poser $v_n=u_n-6$ donne $v_{n+1}=v_n/2$, $v_0=-5$, donc $u_n=6-5/2^n$. La suite croît vers $6$.

\section*{Exercice 10.5 -- Indétermination}
$(2/n)/(1/n)=2$, tandis que $(1/n)/(1/n^2)=n\to+\infty$. Les morceaux ont les mêmes limites nulles, ce qui démontre qu'une forme $0/0$ ne détermine aucune réponse.

\section*{Exercice 10.6 -- Somme infinie}
C'est la limite de $S_n=\sum_{k=1}^n3^{-k}=(1/3)(1-3^{-n})/(1-1/3)$. Comme $3^{-n}\to0$, la limite vaut $1/2$. La définition passe d'abord par une somme finie.
